\section{Introduction}
\label{sec:introduction}

The deployment of autonomous multi-agent AI systems into consequential domains — financial operations, infrastructure management, clinical decision support, and content publishing — has outpaced the development of accountability architectures capable of surviving the failure modes these systems exhibit. Current generation multi-agent frameworks treat accountability as a governance convention: a set of policies, access-control lists, and review workflows that are applied after the system has acted. This approach is insufficient. When a network of autonomous agents coordinates to decompose and execute a complex directive, the causal chain between the human principal's intent and the terminal outcome is mediated by multiple machine actors in ways that are opaque, asynchronous, and legally ambiguous. A governance convention that operates at the policy layer cannot guarantee that a human remains in the loop when the system itself has entered a state it cannot resolve. The consequence is a structural accountability gap — one that current frameworks are architecturally incapable of closing.

We begin with a simple but forceful observation: accountability is not a property that can be declared; it is a property that must be architecturally guaranteed. A system that has a designated human reviewer for its outputs has a accountability policy, not an accountability architecture. The distinction matters because policies fail under stress — they are circumvented by time pressure, misconfiguration, or the emergent behaviors of complex agent networks. An architecture, by contrast, is enforced by the structure of the system itself and cannot be bypassed without the system undergoing a structural change. The question this paper answers is not whether autonomous agents should be accountable — they must be — but rather what the minimal architectural commitment is that makes accountability structurally不可 bypassed, rather than merely declared.

The failure of current systems to provide this guarantee produces two distinct failure modes, each underappreciated in the literature. The first is the \emph{delegation cascade}: when a human principal issues a directive to an agent network, each agent in the network may further delegate subtasks to sub-agents or adjacent agents. The result is a delegation tree in which the terminal actions are not directly authorized by the human principal, but rather authorized through a chain of machine intermediaries. When something goes wrong in such a system — a sub-agent produces a harmful output, a lateral delegation crosses a semantic boundary, or an exception propagates in an unanticipated direction — the system's recovery options are limited by the depth of the delegation chain. A flat agent network with no mandatory escalation pathway will attempt to absorb the exception internally, because that is the behavior the system is optimized for. The outcome is that consequential exceptions are handled by agents that lack the epistemic and corrective capacity to do so.

The second failure mode is the \emph{semantic orphan}: an agent node that has drifted from its intended purpose but has not triggered any threshold that would cause it to be flagged as anomalous by its monitoring systems. Semantic drift is not a software error — it is a gradual misalignment between the agent's operational behavior and the principal's intent that emerges from the statistical nature of large language model outputs. A system without a mandatory bailout mechanism has no architectural way to detect that a node has become a semantic orphan, because the system has no definition of mission coherence that operates independently of the node's self-reported state. The node continues to operate, continuing to take actions that are technically within its provisioned scope but no longer aligned with the principal's intent, until a human notices — and in consequential domains, the noticing may come too late.

The \bxthree{} Framework addresses both failure modes through a structural commitment: every agent node in the network operates within a defined operational envelope enforced by the \bounds{} layer, is assessed for mission alignment by the \purpose{} layer, and is anchored to a cryptographically sealed Forensic Ledger maintained by the \fact{} layer. These three layers are not configurable by individual agent nodes; they are architectural invariants. The Bailout Protocol is the failover mechanism of the \bounds{} layer — the one pathway that every node in the network must support, regardless of its operational role, autonomy level, or task complexity.

The core insight this paper contributes is that accountability in autonomous multi-agent systems must have an architectural fallback — a state that the system is guaranteed to enter when its own resolution capacity is exceeded, independent of any individual node's configuration or volition. We call this fallback a \emph{mandatory bailout}, and we distinguish it from conventional exception handling by three properties: it is architecturally enforced (not policy-declared), it is guaranteed to terminate at a human (not routed to another machine actor), and it is non-negotiable (the system cannot opt out of it by declaring nominal normalcy). These three properties define the minimal structural guarantee that accountability requires, and they are the properties that current multi-agent frameworks lack.

The protocol operates as follows. When the \bounds{} layer of any node detects a condition that exceeds the node's provisioned resolution capacity — an exception it cannot classify, a state it cannot evaluate, or a corrective action it cannot execute within its envelope — it does not attempt to absorb or route the exception within the agent network. Instead, it enters the bailout state: it suspends the node's autonomous operations, seals the current state to the Forensic Ledger, formats the diagnostic context into a human-readable escalation payload, and routes that payload to the node's designated human anchor. The human anchor is not a fallback reviewer; it is the architectural endpoint of the bailout pathway. The node does not resume operation until the human provides a resolution directive, which is recorded in the Forensic Ledger by the \fact{} layer.

This architecture has three consequences that distinguish it from prior approaches. First, the bailout is triggered by the \bounds{} layer's assessment of its own resolution capacity — not by a threshold set by the node itself — which eliminates the semantic orphan problem by ensuring that a node cannot declare itself nominal when it has exceeded its envelope. Second, the bailout pathway is defined at the architectural level, not the policy level, which means it cannot be circumvented by time pressure, system load, or configuration drift. Third, the Forensic Ledger provides an append-only, cryptographically sealed record of every bailout event, ensuring that the accountability chain between human principal and terminal outcome is always reconstructable, even in the presence of multi-level delegation and asynchronous operation.

The remainder of this paper is organized as follows. Section~\ref{sec:background} establishes the theoretical foundations: the accountability challenges specific to multi-agent AI systems, the role of human-in-the-loop design in consequential domains, and the fail-safe design principles borrowed from safety-critical systems engineering. It situates these three lineages within the \bxthree{} three-layer model — the \purpose{}, \truth{}, and \bounds{} layers operating over the \fact{} Forensic Ledger — and demonstrates how each lineage addresses a distinct failure mode that the Bailout Protocol is designed to resolve. Section~\ref{sec:state-machine} presents the Bailout Protocol as a state machine, defining the states, transitions, and invariants that govern the protocol's operation across all node types. Section~\ref{sec:implementation} describes the protocol's implementation within the AgentOS runtime, including the watchdog mechanism that monitors node resolution capacity, the escalation pathway that terminates at the human anchor, and the forensic sealing procedure that ensures the diagnostic context is preserved. Section~\ref{sec:evaluation} provides empirical evaluation through simulation of cascade failure scenarios in a multi-agent network, demonstrating that the protocol terminates accountability chains that would otherwise propagate through the agent network without human resolution. Section~\ref{sec:related-work} surveys the related literature on multi-agent accountability, human-in-the-loop AI, and fail-safe design. Section~\ref{sec:conclusion} concludes with a discussion of the protocol's limitations, the conditions under which it is applicable, and the open questions for future research.